\section{The crossing at every \texorpdfstring{$N$}{N}}\label{app:every-n}

This appendix proves Lemma~\ref{lem:bessel-bounds}, the analytic half of
Theorem~\ref{thm:every-n}, the identity behind the certificates and
Proposition~\ref{prop:crossing-eq}. The letters $U$, $U'$, $W$, $W'$, $r$,
$r'$, $t$, $t'$, $\tau$, $\kappa$, $m$, $B$, $D$, $\beta$, $\Delta$ and
$\delta$ are local to this appendix (each is defined where it first shows up).
We round every decimal constant in the safe direction, namely so that the
inequality it enters stays true.

\subsection{Proof of Lemma~\ref{lem:bessel-bounds}}

For integer $n$, $I_n(z)=\frac1\pi\int_0^\pi e^{z\cos t}\cos(nt)\,dt$
\cite[Eq.~10.32.3]{dlmf}, so
$I_0(z)+I_1(z)=\frac1\pi\int_0^\pi e^{z\cos t}(1+\cos t)\,dt$. Substitute
$s=1-\cos t$, so that $\sin t=\sqrt{s(2-s)}$ and $1+\cos t=2-s$, and then
$s=v/z$. This gives
\begin{equation}\label{eq:E-integral}
\cE(z)=\frac1{\sqrt\pi}\int_0^{2z}e^{-v}v^{-1/2}\sqrt{1-\frac v{2z}}\,dv .
\end{equation}
On $[0,1]$ we have $\sqrt{1-y}=1-\sum_{k\ge1}c_ky^k$ with every $c_k>0$,
$c_1=\frac12$ and $c_2=\frac18$, and $\sum_kc_k=1$ by monotone convergence as
$y\uparrow1$. So $\sum_{k\ge3}c_k=\frac38$, and $y^k\le y^3$ gives
$\sum_{k\ge3}c_ky^k\le\frac38y^3$. Hence, for $0\le y\le1$,
\[
g_-(y):=1-\frac y2-\frac{y^2}8-\frac{3y^3}8\le\sqrt{1-y}\le1-\frac y2-\frac{y^2}8
\le g_+(y):=1-\frac y2-\frac{y^2}8+\frac{y^3}8 .
\]
The cubic $g_-$ decreases and $g_-(1)=0$, so $g_-\le0$ on $[1,\infty)$. Also
$8g_+(y)=(y-2)^2(y+1)+2(y-1)^2+2>0$ for $y\ge-1$. Put $y=v/(2z)$ in
\eqref{eq:E-integral}, replace $\sqrt{1-y}$ by $g_\pm$ and extend the integral
to $\infty$. The added piece is at most $0$ with $g_-$ and at least $0$ with
$g_+$, so the lower bound survives with $g_-$ and the upper with $g_+$. Since
$\frac1{\sqrt\pi}\int_0^\infty e^{-v}v^{k-1/2}\,dv=\Gamma(k+\frac12)/\Gamma(\frac12)$
equals $1,\frac12,\frac34,\frac{15}8$ for $k=0,1,2,3$, the two bounds integrate
to the two sides of the lemma. The bound $\cE\le1$ follows in the same way from
$\sqrt{1-y}\le1$. Finally $-\frac3{128z^2}+\frac{15}{512z^3}\le0$ exactly when
$z\ge\frac54$.

\subsection{Proof of Theorem~\ref{thm:every-n} for \texorpdfstring{$N\ge175$}{N >= 175}}

\emph{Step 1. Bounds that hold for every $\zeta\ge4$.} Fix $N\ge175$ and write
$\zeta=\zeta_N$, $U=\zeta+\frac1{8\zeta}$, $r=U/N$ and
$n(\zeta)=\sqrt{8\zeta/\pi}\,e^\zeta$. Then $n(\zeta_N)=N$ by
\eqref{eq:zetaN}, and $\zeta\ge4$ since $n(4)<175$. For $k=1,2,3$ the function
$U^k/n(\zeta)$ decreases on $[4,\infty)$, because its logarithmic derivative is
$k(1-\frac1{8\zeta^2})/U-1-\frac1{2\zeta}<\frac kU-1<0$. The same holds for
$U^4/n$, since $U>4$, and for $\zeta/n$ and $\zeta^2/n$. The values at
$\zeta=4$ give, for every $\zeta\ge4$,
\[
r\le0.023135,\quad\frac{U^2}N\le0.093262,\quad\frac{U^3}N\le0.375960,\quad
\frac{U^4}N\le1.51559,\quad\frac{\zeta(\zeta+1)}N\le0.114776 .
\]

\emph{Step 2. The upper bound.} Let $W=U/(1-r)$, so that $NW/(N+W)=U$, and let
$\delta=W(W+2)/(2N)$. By Lemma~\ref{lem:bracket}(b) and \eqref{eq:bracket-log}
it suffices to show that $G_U:=\psi(W)-\psi(\zeta)+\log \cE(W)+\log(1-\delta)>0$.
We bound the three parts from below. First,
$\psi(W)-\psi(\zeta)=(W-U)+(U-\zeta)+\frac12\log\frac U\zeta+\frac12\log\frac WU$,
where $W-U=\frac{Ur}{1-r}$, $U-\zeta=\frac1{8\zeta}$,
$\log\frac U\zeta=\log(1+\frac1{8\zeta^2})\ge\frac1{8\zeta^2}-\frac1{128\zeta^4}$
and $\log\frac WU=-\log(1-r)\ge r$. Second, Lemma~\ref{lem:bessel-bounds} gives
$\cE(W)\ge1-t$ with $t=\frac1{8W}+\frac3{128W^2}+\frac{45}{512W^3}$. Since
$W\ge U\ge\zeta$ and
$\frac1{8U}=\frac1{8\zeta}-\frac1{8\zeta(8\zeta^2+1)}\le\frac1{8\zeta}-\frac1{64\zeta^3}+\frac1{512\zeta^5}$,
we get
$t\le\tau:=\frac1{8\zeta}+\frac3{128\zeta^2}+\frac{37}{512\zeta^3}+\frac1{512\zeta^5}$
and $\log \cE(W)\ge\log(1-\tau)\ge-\tau-\frac{\tau^2}{2(1-\tau)}$. Third,
$\delta=\frac{Ur}{2(1-r)^2}+\frac r{1-r}$ and
$\log(1-\delta)\ge-\delta-\frac{\delta^2}{2(1-\delta)}$. Adding the three parts
gives $G_U\ge B+D$, where $B$ collects the terms in $\zeta$ alone and $D$ the
terms with a factor $r$,
\[
B=\frac5{128\zeta^2}-\frac{37}{512\zeta^3}-\frac1{256\zeta^4}-\frac1{512\zeta^5}-\frac{\tau^2}{2(1-\tau)},\qquad
D=\frac r{2(1-r)}\Bigl(\frac{U(1-2r)}{1-r}-1-r\Bigr)-\frac{\delta^2}{2(1-\delta)} .
\]
Both $\zeta\tau$ and $\tau$ decrease in $\zeta$, so $\zeta^2B$ increases. At
$\zeta=4$ we have $\zeta\tau=0.135384$, $\tau=0.033846$ and
$\zeta^2B=0.011236$. Hence $B>0$. In $D$ the bracket is at least
$4.03125(1-0.023683)-1.023135>2.912$, so the first term of $D$ is at least
$1.456\,r$. Write $\delta=rm$ with $m=\frac U{2(1-r)^2}+\frac1{1-r}$. Then
$\delta\le0.07255$ and
$rm^2=\frac{U^3}{4N(1-r)^4}+\frac{U^2}{N(1-r)^3}+\frac r{(1-r)^2}\le0.2276$, so
$\frac{\delta^2}{2(1-\delta)}=r\cdot\frac{rm^2}{2(1-\delta)}\le0.1228\,r$.
Hence $D\ge1.33\,r>0$. So $G_U>0$ and $z_N<U$.

\emph{Step 3. The lower bound.} Let $\varepsilon=\varepsilon_N$,
$U'=U-\varepsilon$, $r'=U'/N$ and $W'=U'/(1-r')$, so that
$NW'/(N+W')=U'$. By Step~1, $\varepsilon\le0.114776+\frac1{256}<0.1187$, so
$U'\ge\zeta-0.1187\ge3.88$ and $W'\ge U'$. By Lemma~\ref{lem:bracket}(a)
and \eqref{eq:bracket-log} it suffices to show that
$G_{U'}:=\psi(W')-\psi(\zeta)+\log \cE(W')<0$. We bound the two parts from above.
First,
$\psi(W')-\psi(\zeta)=(U'-\zeta)+\frac{U' r'}{1-r'}+\frac12\log\frac{U'}\zeta-\frac12\log(1-r')$,
where $U'-\zeta=\frac1{8\zeta}-\varepsilon$,
$\frac12\log\frac{U'}\zeta\le\frac12(\frac{U'}\zeta-1)=\frac1{16\zeta^2}-\frac\varepsilon{2\zeta}$
and $-\frac12\log(1-r')\le\frac{r'}{2(1-r')}$. The two terms in $r'$ add up to
$\frac{U'^2+U'/2}{N(1-r')}$. Second, Lemma~\ref{lem:bessel-bounds} gives
$\cE(W')\le1-t'$ with
$t'=\frac1{8W'}+\frac3{128W'^2}-\frac{15}{512W'^3}\in(0,1)$, so
$\log \cE(W')\le-t'$. Here $\frac1{8W'}=\frac1{8U'}-\frac1{8N}$. Write
$U'=\zeta(1+\kappa)$ with $\kappa=\frac1{8\zeta^2}-\frac\varepsilon\zeta$.
Then $\frac1{U'}\ge\frac{1-\kappa}\zeta$, so
$-\frac1{8W'}\le-\frac1{8\zeta}+\frac1{64\zeta^3}-\frac\varepsilon{8\zeta^2}+\frac1{8N}$.
Next, $x\mapsto x/(1-x/N)$ increases and $U'\le U$, so $W'\le U/(1-r)$ and
$\frac1{W'^2}\ge\frac{(1-r)^2}{\zeta^2(1+\frac1{8\zeta^2})^2}\ge\frac{(1-2r)(1-\frac1{4\zeta^2})}{\zeta^2}$.
Finally $\frac{15}{512W'^3}\le\frac{15}{512U'^3}$. Adding, and using
$\varepsilon(1+\frac1{2\zeta})=\frac{\zeta^2+\frac32\zeta+\frac12}N+\frac1{16\zeta^2}+\frac1{32\zeta^3}$,
we get $G_{U'}\le\beta+\Delta-\frac\varepsilon{8\zeta^2}$ with
\[
\beta=-\frac3{128\zeta^2}(1-2r)\Bigl(1-\frac1{4\zeta^2}\Bigr)+\frac{15}{512U'^3}-\frac1{64\zeta^3},\qquad
\Delta=\frac1N\Bigl(\frac{U'^2+U'/2}{1-r'}+\frac18-\zeta^2-\frac32\zeta-\frac12\Bigr).
\]
For $\beta<0$ it suffices that
$\zeta^2/U'^3\le\frac45(1-2r)(1-\frac1{4\zeta^2})$. The left side is at most
$\zeta^2/(\zeta-0.12)^3\le16/3.88^3<0.274$, since this function decreases for
$\zeta>0.12$ (its logarithmic derivative is $\frac2\zeta-\frac3{\zeta-0.12}<0$
there). The right side is at least $0.8\cdot0.9537\cdot0.9843>0.75$. For
$\Delta<0$, note that $U'\le U$ and $r'\le r$, so
\[
\frac{U'^2+U'/2}{1-r'}\le U^2+\frac U2+\frac{U^3/N+U^2/(2N)}{1-r}
\le\zeta^2+\frac\zeta2+0.2667+0.4326,
\]
using $U^2+\frac U2=\zeta^2+\frac\zeta2+\frac14+\frac1{16\zeta}+\frac1{64\zeta^2}$.
This is less than $\zeta^2+\frac32\zeta+\frac38$ because $\zeta\ge4$. Hence
$G_{U'}<0$ and $z_N>U'$.

The constants were recomputed in 40-digit arithmetic, and as a check outside
the proof we evaluated $\zeta^2B$, $D$, $\beta$ and $\Delta$ on
$\zeta\in[4,60]$ with step $0.001$.

\subsection{The certificates}\label{app:certificates}

Let $c,d>0$ be integers. Multiplying \eqref{eq:SN} by $d^{N-1}$ gives
\[
d^{N-1}S_N(c/d)=\sum_jW(N,m_N,j)\,c^jd^{N-1-j},
\]
an integer. By Theorem~\ref{thm:pmf} it is the sum of
$d^{\rep(w)}c^{\chg(w)}$ over the words $w$ in bin $m_N$, that is, the
numerator $n_N(m_N)$ of Remark~\ref{rem:tables} with $(\alpha,\beta)=(d,c)$.
The recursion \eqref{eq:dp} with the weights $d,c$ in place of $p,q$ and
initial values $1$ computes the same sums. So $S_N(c/d)<1$ exactly when this
integer is less than $d^{N-1}$. This is the identity used in Step~4 of the
proof of Theorem~\ref{thm:every-n}.

\subsection{Proof of Proposition~\ref{prop:crossing-eq}}

\emph{Step 1. An exact formula for $R_N$.} Write $w=w_N$, $z=z_N$,
$\Lambda_N=2\Phi(w)/N$ and $\delta=w(w+2)/(2N)$. At $w=w_N$ the display in the
proof of Lemma~\ref{lem:bracket} gives $1\le\Lambda_N\le1/(1-\delta)$ when
$\delta<1$, and \eqref{eq:bracket-log} gives
$\psi(w)+\log \cE(w)=L+c_0+\log\Lambda_N$. Since $w-z=wz/N$, $w/z=1+w/N$ and
$\frac1{8z}=\frac1{8w}+\frac1{8N}$, the remainder
$R_N=\psi(z)-L-c_0-\frac1{8z}$ is
\begin{equation}\label{eq:RN-formula}
R_N=\log\Lambda_N-\Bigl(\log \cE(w)+\frac1{8w}\Bigr)-\frac1{8N}-\frac{wz}N-\frac12\log\Bigl(1+\frac wN\Bigr).
\end{equation}

\emph{Step 2. The lower bound for $N\ge175$.} By the proof of
Theorem~\ref{thm:every-n}, $3.88\le z\le w\le U/(1-r)$. Note that
$\log\Lambda_N\ge0$ and $-\log \cE(w)\ge-\log(1-\frac1{8w})\ge\frac1{8w}$ by
Lemma~\ref{lem:bessel-bounds}, since $w\ge\frac54$. So
$R_N\ge-\frac{wz+w/2+1/8}N=-\frac{z^2+z/2}{N-z}-\frac1{8N}$. This is at least
$-\frac{z^2+z}N$ as soon as $N\ge\frac{2z^2(z+1)}{z-1/4}$. The right side
increases in $z$ and $z\le U$, so it suffices that
$U^2/N\le\frac{U-1/4}{2(U+1)}$. The left side is at most $0.0933$ and the
right side at least $0.3757$.

\emph{Step 3. The upper bound for $N\ge175$.} Let
$t=\frac1{8w}+\frac3{128w^2}+\frac{45}{512w^3}\le\frac{1.09504}{8w}$, which
holds as $w\ge3.88$. Lemma~\ref{lem:bessel-bounds} and
$-\log(1-t)\le t+\frac{t^2}{2(1-t)}$ give
\[
-\log \cE(w)-\frac1{8w}\le\frac3{128w^2}+\frac{45}{512w^3}+\frac{t^2}{2(1-t)}\le\frac{0.0559}{w^2}
\le\frac{0.0559}{z^2}.
\]
Using $\log\Lambda_N\le\delta+\frac{\delta^2}{2(1-\delta)}$,
$\log(1+x)\ge x-\frac{x^2}2$ and $wz\ge w^2-w^3/N$, the rest of
\eqref{eq:RN-formula} is at most
$-\frac{(w-1/2)^2}{2N}+\frac{w^3+w^2/4}{N^2}+\frac{\delta^2}{2(1-\delta)}$. By
the bounds on $r$ and $U^k/N$ in Step~1 of the previous proof, with
$w\le U/(1-r)$, we have $\frac{w^3+w^2/4}N\le0.4278$ and
$N\frac{\delta^2}{2(1-\delta)}=\frac{(w^4+4w^3+4w^2)/(4N)}{2(1-\delta)}\le0.4945$,
while $\frac{(w-1/2)^2}2\ge5.71$. So the rest is negative and
$R_N\le0.0559/z^2<\frac1{16z^2}$.

\emph{Step 4. The form in $L$ for $N\ge175$.} We have
$z<U=L+c_0-\frac12\log\zeta+\frac1{8\zeta}<L$. Also
$z>\zeta-0.1187\ge\nu_N-0.1187\ge L/2$, because $L-\log L\ge1.18$ for
$L\ge\log175$. Then $\frac1{16z^2}\le\frac1{4L^2}$ and
$\frac{z^2+z}N\le\frac{L^2+L}N$.

\emph{Step 5. $N\le174$.} The function $\phi(z)=\psi(z)-L-c_0-\frac1{8z}$
increases in $z$, and $R_N=\phi(z_N)$, while $\frac1{16z^2}$ and
$-\frac{z^2+z}N$ decrease. So it suffices to check
$\phi(z(c'/d))\le\frac1{16z(c'/d)^2}$ and
$\phi(z(c/d))\ge-\frac{z(c/d)^2+z(c/d)}N$ with the certificates, together with
$z(c/d)\ge L/2$ and, for $N\ge3$, $z(c'/d)\le L$. All hold, and the smallest
gaps are $0.0373$ (upper) and $0.0813$ (lower), both at $N=174$. For $N=2$ we
have $S_2(u)=2u$, so $u_2=\frac12$ and $z_2=\frac23$, and $\frac23<\log2$ is
equivalent to $e^2<8$, which holds since $e^2<2.72^2=7.3984$.
